\documentclass[11pt]{article}
\usepackage[T1]{fontenc}
\usepackage[utf8]{inputenc}
\usepackage[english]{babel}
\usepackage[margin=1in]{geometry}
\usepackage{amsmath,amssymb}
\usepackage{enumitem}
\usepackage{microtype}

\ifdefined\pdfcatalog
  \pdfcatalog{/Lang (en-US)}
\fi

\setlength{\parindent}{0pt}
\setlength{\parskip}{0.55em}
\setlength{\emergencystretch}{2em}
\setlist{nosep,leftmargin=*}

\newcommand{\findinglabel}[1]{\textbf{#1}}
\newcommand{\classification}[1]{\textit{Classification: #1.}}

\title{Technical Review of ``Early-Stopping Activation Steering for Factual Preference Alignment in Vietnamese Medical Language Models''}
\author{Manuscript review}
\date{August 25, 2026}

\begin{document}
\selectlanguage{english}
\maketitle

\section*{Overall assessment}

The latest revision makes several useful local corrections.  The Abstract now acknowledges the approximately 30.5\% relative Rep-4 increase, the principal RAG paragraph gives the correct 97.4\% latency increase and 7.13~s delay, and an empirical randomization value is added to the placebo summary.  The natural-completion table and bounded 200-token McNemar result remain the clearest empirical components.

The manuscript is still not submission-ready because the corrections are not propagated consistently.  The activation Results retain mathematically false differences, the Abstract still claims a $2.67\times$ context expansion contradicted by its own prompt-token table, and the Conclusion retains the obsolete 6.38~s latency value.  Experimental Setup and Results use different RAG corpus sizes, while Discussion still reports an obsolete natural-completion effect.  Beyond these definite errors, activation norm is still interpreted as stabilization without a controlled mechanistic test, RefPref remains a clinical proxy, and the RAG, placebo, latency, dataset, vector, and implementation analyses remain incomplete.  Major revision is required.

\section*{Major technical findings}

\subsection*{1. The activation section still contains contradictory arithmetic and conclusions}

\classification{Definite quantitative errors}

\findinglabel{Location.} Abstract; contribution~3; Empirical Representation Saturation Diagnostics.

\findinglabel{Problem.} The Abstract correctly states that continuous steering is 1.97 units above baseline at $t=50$ and linear decay is 3.01 units below baseline.  The Results paragraph still states that 56.18 versus 54.21 is a $+2.93$ distortion and that 51.20 is within 0.19 units of 54.21.  Direct calculation gives
\[
56.18-54.21=1.97,
\qquad
51.20-54.21=-3.01.
\]
The Abstract then describes the 3.01-unit undershoot as ``stabilizing magnitude inflation,'' although no stabilization criterion or acceptable range is defined.  The contribution list reports ``linear decay stabilization (51.20)'' without noting that this value is farther from the displayed baseline than the continuous-steering value is.  The hard-cutoff change from 56.28 to 52.75 still lacks time indices.

\findinglabel{Why it matters.} The activation evidence is presented as mechanistic validation.  Two incompatible interpretations of the same values prevent readers from determining whether the diagnostic supports convergence, undershoot, or neither.

\findinglabel{Suggested fix.} Generate the Abstract, contribution, Results, and Discussion from one activation-summary artifact.  Report the full trajectory, time indices, mean, dispersion, confidence interval, and per-step sample count.  Remove ``stabilization'' unless a prespecified baseline range or functional criterion establishes it.

\subsection*{2. Norm elevation and later undershoot do not establish a saturation-prevention mechanism}

\classification{Unsupported mechanism claim and implementation ambiguity}

\findinglabel{Location.} Abstract; Introduction; Early-Stopping Steering Hook \& Saturation Diagnostics; activation Results.

\findinglabel{Problem.} A scalar $\ell_2$ norm does not establish saturation, manifold departure, or recovery of a representation.  Saturation requires an operational property such as clipping, nonlinear loss of responsiveness, directional collapse, or a measured behavioral degradation associated with the state.  A scalar radius cannot determine whether an activation is outside a high-dimensional manifold.  Turning the intervention off after $K=16$ also makes later removal of the direct additive term expected by construction, not evidence that decay repairs an internal failure.

The manuscript does not specify whether the diagnostic consistently measures pre-hook $h_l^{(t)}$ or post-hook $\hat h_l^{(t)}$, how norms are aggregated, which tensor positions are altered, how sequences ending before $t=50$ or $t=100$ are handled, or whether schedules are evaluated on identical token contexts.  In free generation, different previous tokens confound activation comparisons.

\findinglabel{Why it matters.} The central novelty is motivated as controlling a representation-level failure, but the current diagnostic only provides under-specified scalar summaries.

\findinglabel{Suggested fix.} Define the failure operationally.  Add fixed-token replay or teacher forcing, pre/post-hook trajectories, survivor counts, confidence bands, projection onto $v_{\text{steer}}$, cosine drift, distributional distance, and response to an additional perturbation.  Relate these diagnostics to repetition and factual errors.  Use ``norm elevation/undershoot'' rather than ``saturation,'' ``manifold departure,'' or ``stabilization'' until those stronger claims are demonstrated.

\subsection*{3. RAG context size and systems summaries remain internally inconsistent}

\classification{Definite data-definition and section-to-section errors}

\findinglabel{Location.} Abstract; RAG Results; Table~\texttt{tab:baselines}; Conclusion.

\findinglabel{Problem.} The main RAG paragraph now correctly reports
\[
14.45-7.32=7.13\ \text{s},
\qquad
\left(\frac{14.45}{7.32}-1\right)100\%\approx97.4\%.
\]
The Conclusion still reports ``+6.38s latency.''  The Abstract still claims a $2.67\times$ prompt-context expansion, while Table~\texttt{tab:baselines} reports 118.0 prompt tokens for baseline and 75.4 for RAG:
\[
75.4/118.0\approx0.639,
\]
which is a 36.1\% reduction, not an expansion.  The RAG paragraph says 75.4 tokens incorporate retrieved passages and inflate KV-cache requirements, which also conflicts with the displayed baseline count.  If 75.4 denotes retrieved passage tokens only, the table header ``Prompt Tokens'' is incorrect and the total RAG prompt is not reported.

\findinglabel{Why it matters.} The efficiency and memory comparison is a headline result, but the manuscript uses incompatible quantities and stale values.

\findinglabel{Suggested fix.} Define the token-count boundary, including system instructions, user query, and retrieved passages.  Recompute the total prompt length and all derived ratios from raw logs.  Propagate 7.13~s and 97.4\% consistently if 7.32/14.45 are authoritative.  Remove the $2.67\times$ and 6.38~s claims unless supported by a different documented measurement.

\subsection*{4. The BM25 experiment remains non-reproducible and too weak for a general RAG conclusion}

\classification{Definite corpus inconsistency and likely weak baseline}

\findinglabel{Location.} Abstract; Experimental Setup; RAG Results; Table~\texttt{tab:baselines}.

\findinglabel{Problem.} Experimental Setup specifies 4,495 passages, while Results specifies 14,576.  The Abstract calls 19.20\% Recall@1 ``retrieval noise,'' although it is a success rate; the corresponding Top-1 miss rate is 80.80\%.  Passage construction, Vietnamese tokenization, query formulation, BM25 implementation and parameters, relevance labels, and retrieval-failure handling are not defined.  Oracle Gold-Context RAG is represented only by BERTScore 0.8002, with no RefPref, prompt length, latency, generation configuration, or paired comparison.  No paired interval/test supports the claimed 8.60-point superiority.

\findinglabel{Why it matters.} The corpus mismatch blocks reproduction, and an under-specified low-recall retriever cannot support a broad conclusion about RAG.

\findinglabel{Suggested fix.} Use one versioned corpus and publish its passage IDs and construction procedure.  Report miss rate separately from recall, category-wise Recall@$k$, retrieval-conditioned generation results, complete oracle results, and stronger lexical/dense retrievers.  Compare all methods on identical records using paired uncertainty.  Describe the result as one BM25 configuration.

\subsection*{5. Natural-completion reporting still contains obsolete claims and lacks paired inference}

\classification{Definite contradiction and statistical omission}

\findinglabel{Location.} contribution~2; Primary Natural Completion Evaluation; Discussion; Conclusion.

\findinglabel{Problem.} The table reports baseline 65.40\%, linear decay 67.80\%, continuous steering 68.60\%, and hard cutoff 70.60\%.  Discussion still reports an obsolete $+1.00$-point nonsignificant trend.  Contribution~2 and the first Results paragraph state 100.00\% EOS, whereas continuous steering and linear decay are 99.80\%.  A 99.80\% rate means one of 500 outputs did not emit EOS under the stated criterion.  No paired transition tables, CIs, exact tests, or multiplicity policy are reported for the three schedule-versus-baseline comparisons.

\findinglabel{Why it matters.} Marginal percentages do not establish significance or schedule ranking, and stale text changes both the effect size and termination claim.

\findinglabel{Suggested fix.} Replace the obsolete Discussion and EOS statements.  Report the complete paired tables, effect estimates and CIs for RefPref/BERTScore/Rep-4, exact test definitions, multiple-comparison handling, and the non-EOS finish reason.

\subsection*{6. The natural-completion experiment demonstrates a metric trade-off, not uniform quality improvement}

\classification{Unsupported comparative language}

\findinglabel{Location.} Abstract; natural-completion Results; Discussion; Conclusion.

\findinglabel{Problem.} All steering schedules increase RefPref but decrease reference BERTScore and increase Rep-4.  Hard cutoff raises Rep-4 from 4.10\% to 5.35\%, a 30.5\% relative increase; linear decay raises it by approximately 31.0\%.  Hard cutoff is best for RefPref, baseline is best for BERTScore and Rep-4, and linear decay does not dominate continuous steering.  The Conclusion still says linear decay yields ``consistent reference-preference improvements'' while the strongest natural-completion RefPref result belongs to hard cutoff and no paired uncertainty is given.

\findinglabel{Why it matters.} ``Improvement'' can be mistaken for better overall or clinical quality when the measured effects move in opposing directions.

\findinglabel{Suggested fix.} Report paired metric differences with CIs and define the primary endpoint.  If the intended claim is improved RefPref without unacceptable degradation, prespecify noninferiority margins for BERTScore and Rep-4.  Select the primary schedule on validation data and name the metric whenever using comparative language.

\subsection*{7. The RefPref endpoint does not establish clinical correctness or hallucination mitigation}

\classification{Unsupported interpretation of a disclosed proxy}

\findinglabel{Location.} Abstract; Introduction; contributions; Results; Conclusion; keywords.

\findinglabel{Problem.} The Methods correctly calls RefPref a semantic proxy, but the paper still describes the vector as ``truthfulness,'' uses ``factual alignment'' and ``accuracy,'' and retains hallucination-mitigation framing.  Relative similarity to $y^+$ rather than one synthetic $y^-$ cannot determine whether a response contains a wrong dose, unit, negation, contraindication, or interaction.  The natural results show RefPref increasing while reference BERTScore decreases, further demonstrating the endpoint's non-equivalence to direct correctness.

\findinglabel{Why it matters.} Clinical safety and deployment claims require factual and harmful-error labels, not only embedding similarity.

\findinglabel{Suggested fix.} Add blinded clinician-adjudicated correctness and harmful-error endpoints, category stratification, inter-rater agreement, adjudication, and RefPref validation.  Otherwise use ``reference-preference rate'' consistently and remove clinical-accuracy or hallucination-mitigation claims as measured outcomes.

\subsection*{8. The placebo randomization result is minimally resolved and incompletely specified}

\classification{Likely statistical overstatement}

\findinglabel{Location.} Abstract; contribution~4; directional-control Methods; placebo Results; Conclusion.

\findinglabel{Problem.} The Abstract now reports empirical randomization $p=0.0476$, which is exactly $1/(20+1)$ and therefore the smallest plus-one-corrected value available with 20 random directions when none exceeds the target.  The manuscript does not state that the test is one-sided, define the statistic or plus-one convention, or provide the 20 seeds and individual scores.  The body and Conclusion still emphasize ``$+4.06\sigma$,'' which is only a standardized distance, not a calibrated Gaussian-tail test.  BERTScore dispersion and hierarchical uncertainty over questions and directions remain absent.  Positive and negative steering are markedly asymmetric.

\findinglabel{Why it matters.} The experiment provides limited direction-specific evidence, but its inferential resolution is borderline and cannot prove that the vector encodes factuality.

\findinglabel{Suggested fix.} Define the empirical test, alternative hypothesis, exchangeability assumption, and correction; release seeds and individual outputs; increase the number of random directions; and use hierarchical bootstrap uncertainty.  Replace the $4.06\sigma$ language with the actual randomization result and discuss asymmetry.

\subsection*{9. Systems timing, ablation rankings, and the matched-dose definition remain incomplete}

\classification{Implementation ambiguity and definite formula inconsistency}

\findinglabel{Location.} dose equations; Experimental Setup; baseline and ablation tables; Discussion.

\findinglabel{Problem.} Every non-RAG condition remains exactly 7.32~s, without repetitions, variance, warm-up, synchronization, batch size, GPU count, timing boundaries, or peak-memory measurements.  ``No measurable overhead'' is not established by an equivalence interval.  Ablation differences of 0.0005--0.0018 BERTScore lack paired CIs and multiplicity handling, while temporal placement and $D_2$ remain unmatched.  Finally,
\[
\alpha_{\mathrm{match}}=D_{\mathrm{decay}}/T
=\frac{K+1}{2T}\alpha_0
\]
equates a nonnegative absolute dose with a signed expression when $\alpha_0<0$.

\findinglabel{Why it matters.} Efficiency is unmeasured statistically, schedule ranks may be noise, and the dose formula is not valid over the stated negative-steering domain.

\findinglabel{Suggested fix.} Run repeated component-wise timing/memory measurements with an equivalence margin.  Report paired ablation CIs and exploratory multiplicity handling.  Define $\alpha_{\mathrm{match}}=\operatorname{sgn}(\alpha_0)D_1/T$ and $\sum_t|\alpha_{\mathrm{match}}|=D_1$, or restrict the control to positive $\alpha_0$.

\subsection*{10. Dataset, vector, layer, metric, and hook details remain insufficient for reproduction}

\classification{Reproducibility limitation and likely leakage risk}

\findinglabel{Location.} benchmark construction; vector estimation; layer sweep; metric definitions; Experimental Setup.

\findinglabel{Problem.} ``Expert-structured'' is undefined, with no annotation protocol, counter-answer rules, clinical review, deduplication, source IDs, licensing, or source/template-grouped splits.  The 500-item selection rule and seed are missing.  The direction is extracted at the final token of a completed answer and injected into early states without prompt-only, initial-token, pooled, or independent-error controls.  The layer experiment is only a steering sweep range but is called subspace probing and used to infer distributed truthfulness.  Rep-4, Vietnamese ROUGE tokenization, complete BERTScore settings, statistical calls, hook module/indexing, tensor positions, KV-cache behavior, versions, seeds, and GPU count are unspecified.

\findinglabel{Why it matters.} Leakage and implementation choices may materially affect every result and prevent replication.

\findinglabel{Suggested fix.} Add a dataset card and grouped split manifests, publish the evaluation seed, document expert validation, test alternative vector locations/layers, provide complete metric and statistical configurations, and release executable hook pseudocode and per-question outputs.

\section*{Equation, notation, sign, branch, and dimensional audit}

The direction $\mu_l^+-\mu_l^-$ and intervention $+\alpha(t)v_{\text{steer}}^{(l)}$ are sign-consistent.  Unit normalization makes the vector dimensionally compatible with the hidden state when $\alpha(t)$ is an activation-scale coefficient.  The linear schedule and positive-$\alpha_0$ values of $D_{\text{continuous}}$, $D_{\text{cutoff}}$, and $D_{\text{decay}}=(K+1)|\alpha_0|/2$ are arithmetically correct.  No inverse-trigonometric functions, coordinate transformations, or branch/quadrant choices occur.

The concrete issues are the signed/absolute mismatch in $\alpha_{\text{match}}$, ambiguity between pre-hook $h_l^{(t)}$ and post-hook $\hat h_l^{(t)}$, and overload of $N$ for training examples, categories, prompts, and random directions.  The discrete token range should be $t\in\{1,\ldots,100\}$ rather than $t\in[1,100]$.  Terminology drifts among ``truthfulness,'' ``factual preference,'' and RefPref, and among ``norm elevation,'' ``stabilization,'' and the remaining ``saturation diagnostics'' heading.

\section*{Meaningful editorial, compilation, and reference findings}

\subsection*{The source compiles but the experimental list is malformed and tables overflow}

\findinglabel{Location.} Experimental Setup; Tables~\texttt{tab:long\_gen} and \texttt{tab:mcnemar}; clean pdfLaTeX log.

\findinglabel{Problem.} The heading is now valid LaTeX, but the two evaluation modes are still written as plain numbered text rather than an \texttt{enumerate} environment.  After two pdfLaTeX passes, the source builds to six pages; the log reports overfull boxes of approximately 61.44~pt for the natural-completion table and 5.54~pt for the McNemar table.

\findinglabel{Why it matters.} The natural table can intrude far into the adjacent IEEE column, and the list is not structurally typeset.

\findinglabel{Suggested fix.} Use an \texttt{enumerate} environment and redesign or resize both tables while preserving legibility.  Rebuild twice and inspect the final log and PDF.

\subsection*{Citation fit remains uneven}

\findinglabel{Location.} Introduction; bibliography entries \texttt{ref10}, \texttt{ref11}, and \texttt{ref12}.

\findinglabel{Problem.} \texttt{ref10} is the LoRA paper rather than direct evidence for broad SFT hallucination mitigation.  \texttt{ref11} addresses irrelevant retrieved context but does not establish that internal activation pathways cause evidence ignoring.  \texttt{ref12} concerns catastrophic forgetting generally rather than the claimed medical-language-model failure mode.

\findinglabel{Why it matters.} Direct evidence and hypotheses are presented with indistinguishable citation strength.

\findinglabel{Suggested fix.} Add direct sources for hallucination-focused SFT, language-model forgetting after domain adaptation, and evidence use in RAG.  Present the activation explanation as a hypothesis unless directly supported.

\section*{Proofing sweep}

The bundled scan was run on both the source and clean compiled PDF and returned no high-confidence regex hits.  Manual checks identified the following bounded corrections:

\begin{itemize}
    \item \textbf{Contribution~2 and first natural-completion paragraph:} replace the universal 100.00\% EOS statement with the condition-specific 99.80\%--100.00\% range.
    \item \textbf{Activation Results:} remove the stale ``$+2.93$'' and ``within 0.19 units'' statements; state the actual 1.97 and 3.01 differences.
    \item \textbf{RAG Abstract/Conclusion:} remove the stale $2.67\times$ context expansion and 6.38~s delay unless supported by corrected token/time definitions.
    \item \textbf{Retrieval terminology:} call 19.20\% a Recall@1 success rate and 80.80\% the Top-1 miss rate, rather than calling recall itself noise.
    \item \textbf{Experimental Setup:} convert the two plain numbered paragraphs to a LaTeX list.
\end{itemize}

The bibliography was checked for duplicated punctuation, malformed quotation punctuation, and obvious citation-format glitches; no additional mechanical defect was found.

\section*{Items requiring external verification}

\begin{itemize}
    \item Verify all formulary-derived answers and synthetic counter-answers against the official source, current clinical guidance, qualified clinician review, and reuse permissions.
    \item Validate multilingual BERTScore for Vietnamese clinical negation, quantities, units, contraindications, and interaction mechanisms.
    \item Verify novelty against current token-dependent, schedule-based, and dose-controlled activation-steering work.
    \item Verify the description of Qwen2.5-7B-Instruct as an SLM and the privacy/local-hospital claims with explicit definitions, a threat model, and measured resource requirements.
    \item Verify BM25 and stronger Vietnamese retrieval baselines before generalizing the RAG comparison.
\end{itemize}

\section*{Highest-priority revision sequence}

\begin{enumerate}
    \item Propagate one authoritative activation, RAG timing, prompt-token, corpus-size, and natural-completion analysis into every section.
    \item Redesign the activation experiment to test an operational mechanism on controlled contexts with uncertainty.
    \item Add paired natural-completion inference and accurately present the metric trade-off.
    \item Add clinician-adjudicated correctness/harm endpoints or narrow all claims to reference preference.
    \item Strengthen and fully specify BM25/oracle RAG, random-direction inference, and systems timing.
    \item Add ablation uncertainty, correct the matched-dose sign definition, and test vector/layer alternatives.
    \item Complete dataset provenance, leakage controls, implementation details, list/table formatting, and citation fit.
\end{enumerate}

\end{document}